This paper investigates whether retrieval-augmented prompting
improves vision-language model (VLM) accuracy on medical
visual question answering (VQA), using only openly licensed,
non-credentialed data. We build a CLIP-embedded FAISS
retrieval index over the ROCOv2 radiology corpus (approximately
90{,}000 image-caption pairs) and use it to supply retrieved
evidence captions to a pretrained BLIP-2 vision-language model,
comparing this retrieval-augmented condition against the same
model answering with no retrieved evidence. Both conditions
share identical tokenization, decoding parameters, and scoring
code, differing only in prompt content, and are evaluated on the
official VQA-RAD held-out test split (451 examples: 251
closed-ended, 200 open-ended) using BLEU-4, ROUGE-L, and exact
match, together with paired bootstrap significance testing.
Retrieval augmentation improved every reported metric at the
default retrieval depth ($k=5$), with closed-question accuracy
rising from 43.0\% to 47.8\% and overall exact match from
26.8\% to 29.5\%; however, neither improvement was statistically
significant (closed-question exact match: $p=0.1848$, 95\% CI
$[-0.020, 0.116]$; open-question ROUGE-L: $p=0.9642$, 95\% CI
$[-0.032, 0.033]$; $n=451$). A secondary finding from a sweep over
$k \in \{1,3,5,10\}$ showed accuracy peaking at $k=1$ and declining
monotonically as retrieval depth increased, suggesting that excess
retrieved evidence can dilute rather than reinforce a base model's
answer. These results indicate retrieval augmentation shows a
promising but unconfirmed effect on medical VQA at this scale, with
retrieval depth mattering at least as much as the presence of
retrieval itself.
